\section{Discussion}

This benchmark addresses a deliberately narrow practical question: when cell labels are predicted from a fixed frozen-scGPT representation, how does the complete pipeline compare with supervised models trained on the same selected expression variables? The answer was atlas-dependent. Matched expression outperformed scGPT embeddings with every prespecified classifier in TNBC and PBMC, whereas scGPT embeddings were higher with all three classifiers in brain. Donor-cluster intervals supported all six TNBC and PBMC expression advantages, supported the scGPT advantage for brain Random Forest, and did not separate the brain Logistic Regression or XGBoost pipelines. The resulting conclusion is not that either representation is universally superior. It is that frozen scGPT utility depends on tissue, label structure and downstream learner, and that strong matched-expression models remain necessary controls.

The donor-held-out design changes the meaning of the benchmark. Randomly separating cells within an atlas estimates interpolation when donor, sample and batch signatures can occur on both sides of the split. Grouping every donor's cells together instead asks whether a labelled training cohort supports annotation of an unseen donor processed within the same curated resource. The evaluated test sets therefore exclude direct donor overlap, but they do not constitute external-atlas validation. Acquisition protocols, annotation conventions, ancestry composition and study-specific curation remain shared within each resource. The appropriate interpretation is within-atlas unseen-donor performance, not broad biological generalization.

Fairness between representations requires more than supplying the same classifier name. The official scGPT embeddings are unit-normalized and 512-dimensional, while matched expression is sparse, differently scaled and approximately twice as wide. A fixed Logistic Regression penalty can act differently in these spaces. Independent searches over regularization and scaling remove that avoidable asymmetry. The same principle applies to tree learners, whose feature-sampling and depth choices interact with representation geometry. Keeping the classifier fixed for each primary contrast also prevents a favourable head from being chosen after observing outer-test performance.

The SVD-512 control clarifies the role of compression. A training-fitted 512-component expression projection retained most of the uncompressed expression performance, with Logistic Regression differences of only 0.0005--0.0049 across the three atlases. SVD-512 preserved the matched-expression advantage over scGPT in TNBC and PBMC and the scGPT advantage in brain. Thus, the results are not consistent with a simple explanation in which reducing approximately 1,100 selected genes to 512 coordinates necessarily causes the observed losses. At the same time, SVD is not a substitute for scVI, does not model counts generatively and does not remove batch effects. Generative latent-variable baselines remain an important extension.

The brain result is an informative counterexample to a one-directional narrative. Rare macrophage, vascular and precursor classes contributed strongly to Macro F1 despite accounting for a small fraction of cells. scGPT may preserve features useful for separating some rare brain populations, or its pretraining distribution may be more favourable for this tissue. The present design cannot distinguish those mechanisms. Only 12 brain donors were available, and the conditional intervals were correspondingly wider than in TNBC and PBMC. The Random Forest interval nevertheless favoured scGPT, while Logistic Regression and XGBoost remained uncertain at the donor-resampling level. Reporting all three fixed heads prevents one of these outcomes from being elevated post hoc.

The strict scarcity experiment did not reveal a reversal in favour of frozen scGPT embeddings. Matched expression had higher Macro F1 at every tested label budget and in every outer fold even when the gene list was learned only from labelled cells. At the smallest budget, both representations exceeded 0.77 mean Macro F1 across 13 PBMC categories, showing that the supervised heads extracted useful signal from very limited labels. Calibration told a different story: scGPT embeddings sometimes gave lower uncalibrated ECE at larger budgets despite lower discrimination. This divergence reinforces that accuracy, calibration and representation utility are separate properties.

Calibration must be attributed to the complete prediction pipeline. Frozen scGPT supplies a vector, not the class probability used here; the supervised head and its regularization shape confidence. Random Forest was substantially less calibrated than Logistic Regression and XGBoost in TNBC and PBMC for both representations. Temperature scaling improved mean ECE, Brier score and NLL in the TNBC and brain sensitivity analysis, but improvement was not universal at the individual fold level. A scalar temperature cannot repair class confusion or missing biological information, and improved ECE does not guarantee simultaneous improvement in every proper score. Reliability-bin counts, Brier score and NLL should therefore accompany ECE, preferably under external cohort shift.

Class-level results similarly limit any single-number account. In TNBC, the main difference was not a general failure of one representation but a concentration of errors at B-lineage, T-cell and myeloid boundaries. In PBMC, naive and conventional CD8 T cells overlapped with CD4 T and natural-killer categories, with larger errors from scGPT embeddings. Brain errors clustered among biologically related or rare populations, including central-nervous-system macrophages versus microglia and vascular-associated cells versus oligodendrocytes. Portal categories provide reproducible targets, but they are not an error-free biological ground truth. Differences may reflect transcriptional continuity, donor-specific abundance, annotation granularity or label harmonization as well as representation quality.

The pig analysis illustrates why dataset count should not be confused with independent evidence. One donor value cannot support donor-held-out inference, and direct uppercase symbol matching is not a defensible orthology workflow. The high random-split scores show only that the pipelines can interpolate broad categories within this object. A dedicated cross-species evaluation would require an explicit orthologue map, species-aware preprocessing, multiple animals and animal-held-out validation.

The uncertainty analysis improves the biological unit of inference but remains conditional. Donor-cluster resampling preserves within-donor dependence and pairs the representation contrast. It does not refit HVG selection, SVD or models. Repeating the complete nested workflow inside each bootstrap replicate would quantify additional training and selection variability but would be substantially more computationally demanding. The reported intervals therefore show sensitivity to the observed donor composition after fitting. Independent cohorts remain the strongest test of reproducibility.

Several limitations define the evidence boundary. The human benchmark contains three atlases and does not cover the full diversity of tissues, technologies or disease contexts; in particular, a completed donor-held-out TME or human-pancreas analysis is not part of the evidence. SVD-512 Random Forest and XGBoost were not available as complete five-fold pipelines for PBMC and were omitted rather than imputed. Temperature scaling was restricted to TNBC and brain. Only scGPT 0.2.5, one released whole-human checkpoint and one frozen embedding interface were studied. Official end-to-end fine-tuning, newer checkpoints, scVI-like baselines and other foundation models were not tested. Public checkpoint pretraining may overlap source studies, and available metadata do not support a complete cell-level overlap audit. The target labels are broad and differ in granularity among atlases. Finally, the benchmark concerns supervised annotation of known categories rather than novel-state discovery, perturbation response, regulatory inference or multimodal transfer.

Within those boundaries, the study provides a reproducible evaluation template. Foundation-model embeddings should be compared with strong matched inputs, tuned on training groups, assessed under biological holdout and evaluated as complete representation--classifier pipelines. Compression controls and class-level errors are needed to interpret performance differences, while calibration should be reported with multiple probability metrics and a clearly identified calibration cohort. These practices remain useful whether a pretrained representation wins, ties or loses.
